\section{A localized gradient inequality}\label{sec:localized}

We now work over the reals.  The two inputs from the stationary argument
are the analytic normalized family and its polynomial jets in
Proposition~\ref{lin:family}, and the finite-power estimate in
Lemma~\ref{jet:finite-power}.  The latter is used at one fixed Taylor
order.  No convergence of the series in the opening variables is needed
in this section.

If there are no cylindrical opening variables, use $f=F(S_b)$ and
$T_N=S_b$, omit the vacuous finite-power step and the opening-dependent
cutoff, and take $\gamma=1$.  For the otherwise unused order notation
one may take $m=1$ and $N\ge6$.  Keep the actual-graph cutoff and
the compact-source terms in the localized comparison; the subsequent
flow argument is unchanged.

Translate the singular spacetime point to $(0,0)$, and write the original
flow at negative times as $M_t$.  Its rescaling is
\begin{equation}\label{flow:rescaling}
 M(\tau)=e^{\tau/2}M_{-e^{-\tau}},\qquad
 \Phi=\mathbf H+\tfrac12 X^\perp,\qquad
 d(\tau)^2=\int_{M(\tau)}|\Phi|^2e^{-|X|^2/4}\,d\mu .
\end{equation}
We omit the constant Gaussian normalization, consistently with the
stationary sections.  Let $S$ be the time-minus-one section of the
tangent chosen in Proposition~\ref{geo:tangent}.  In particular $S$ is
smooth, properly embedded, of multiplicity one, and has the stated
finite collection of conical and round cylindrical ends.  The same
proposition supplies smooth convergence on compact subsets along the
chosen tangent sequence and a constant $\Lambda$ such that every
rescaled slice under consideration satisfies
\begin{equation}\label{flow:area-ratios}
 \mathcal H^2(M(\tau)\cap B_r(y))\le \Lambda r^2
 \quad(y\in\mathbb R^3,\ r>0).
\end{equation}
All constants below may depend on $S$ and $\Lambda$.

Huisken's monotonicity formula \cite{Huisken1990} gives
\begin{equation}\label{flow:energy}
 \mathcal E(\tau):=F(M(\tau))-F(S)\ge0,\qquad
 \mathcal E'(\tau)=-d(\tau)^2,\qquad \mathcal E(\tau)\longrightarrow0.
\end{equation}
Here the last assertion and the value of the limit require the full
Gaussian integral, rather than only convergence on a compact set.
Indeed, for every fixed $a>0$ and fixed integer $k$, the area-ratio
bound, applied to annuli, gives
\begin{equation}\label{flow:gaussian-tail}
 \int_{M(\tau)\setminus B_R}(1+|X|)^k e^{-a|X|^2/4}\,d\mu
 \le C_{a,k,\xi}e^{-(a-\xi)R^2/4}\quad(0<\xi<a).
\end{equation}
The estimate is uniform in $\tau$, and also holds for the normalized
comparison family below.  It identifies the monotone limit with
$F(S)$ along the chosen tangent sequence, proving \eqref{flow:energy}.

Shrink to a closed parameter neighborhood for $b$ in the family $S_b$,
with $S_0=S$.  The family has uniform bounded geometry in each fixed
derivative order, a positive normal tubular radius, and uniform smooth
end asymptotics.  Its normal shrinker defect is supported in one fixed
compact set $K_0$ and is a smooth compact source with coefficients
$\lambda(b,0)$.  End rotations are incorporated in the parametrizations
of $S_b$; a normal graph over $S_b$ below means physical normal height,
not displacement from $S$ in a fixed unrotated end chart.  Compact
observations are always those fixed in Lemma~\ref{lin:observations}.
Their values determine $(b,x)$ by the fixed analytic coordinate map.

For precision, a complete graph on a ball means that the \emph{entire}
intersection of the surface with that ball is given by the indicated
single graph: no further component or sheet is permitted.  Graph
smallness in the weak $C^1$ tolerance means
$\sup(|u|+|\nabla_{S_b}u|)\le\varepsilon_{\rm weak}$ for its
physical normal height $u$; ``weak'' distinguishes this preliminary
tolerance from the improved bounds, not a weak topology.  Graph
domains can be specified on the reference surface with a fixed
additive margin at their outer boundary.  Replacing an ambient radius
$R$ by $R-C$ only changes a Gaussian estimate below by an arbitrarily
small loss in its quadratic exponent once $R$ is large.

\begin{proposition}[Localized gradient and opening estimates]
\label{flow:localized}
There are $\theta\in(0,1/2)$ and $\gamma>0$ with the following property.
For every sufficiently small $\eta>0$ there are a neighborhood of $S$
on the observation core, a weak normal height and slope tolerance,
and $R_\eta,C_\eta<\infty$ such that the following holds.  Suppose
$M$ obeys \eqref{flow:area-ratios}, its observations lie in that
neighborhood, and $M$ is a complete graph with that weak tolerance over
$S_b$ to radius $R\ge R_\eta$.  Put
\begin{equation}\label{flow:D}
 \mathcal D=\|\Phi_M\|_{L^2_G(M)}
              +e^{-(1-\eta)R^2/8}.
\end{equation}
Then
\begin{equation}\label{flow:localized-conclusions}
 |F(M)-F(S)|^{1-\theta}\le C_\eta\mathcal D,
 \qquad |x|\le C_\eta\mathcal D^\gamma .
\end{equation}
For each fixed compact reference region $K$, transported to $K_b\subset
S_b$ by the family charts, there are $R_{\eta,K}\ge R_\eta$ and
$C_{\eta,K}<\infty$ such that, when $R\ge R_{\eta,K}$, the normal
height $u$ of $M$ satisfies
\begin{equation}\label{flow:compact-height}
 \|u\|_{H^2(K_b)}\le C_{\eta,K}
        (\mathcal D+\mathcal D^\gamma).
\end{equation}
When $\mathcal D\le1$, the right side can be replaced by
$C_{\eta,K}\mathcal D^\gamma$.  Also
\begin{equation}\label{flow:base-multiplier}
 |\lambda(b,0)|\le C_\eta
       (\mathcal D+\mathcal D^\gamma).
\end{equation}
The exponents $\theta,\gamma$ are independent of $\eta$; the
neighborhood, tolerance, minimum radius, and constants need not be.
\end{proposition}

\begin{proof}
Let $m$ be the fixed exponent in Lemma~\ref{jet:finite-power}.  Choose
one integer $N$ for which that lemma applies and
$N-1>4\max\{m,1\}$.  Use the finite action $f_N$, multiplier
$\lambda_N$, and globally cut-off normal Taylor graph $T_N$ from
Lemma~\ref{ref:finite-taylor}.  Its height $v_N$ satisfies
$\|v_N\|_{C^2}\le C_N|x|^{1/2}=o(1)$, including the first two spatial
derivatives needed in the coefficient subtraction.  If there are no
opening variables, $v_N=0$.  The normalized equation, observation,
action, and gradient errors are precisely
\eqref{ref:taylor-equation}--\eqref{ref:taylor-multiplier}.

Let $u$ be the actual normal height.  Its observations define $p=(b,x)$
by $P(u)=\mathsf P(p)$, and its multiplier is zero.  Apply
Lemma~\ref{ref:real-comparison} with $\lambda_U=0$ and weight
$a=1+\delta$ for a small fixed $\delta>0$.  Take the outer reference
radius to be $R-C$, inside the complete graph domain, and the cutoff
fraction $\vartheta_{\rm cut}<1$ sufficiently close to one.  We continue
to write $w=\chi_R(u-v_N)$ for this cutoff difference.  Choose $R_\eta$
large enough that $\chi_R=1$ on the compact source and observation
supports.  For each prescribed compact reference region $K$, increase
$R_{\eta,K}$ so that $\chi_R=1$ on $K_b$ whenever $R\ge R_{\eta,K}$.
The reference lemma retains the actual principal coefficients on the
height difference and uses the displayed $C^2$ smallness only on the
Taylor height; no actual Hessian bound is part of this application.

The actual residual is the scalar pullback of $\Phi_M$.  Its
$L^2_{1+\delta}$ norm is at most $C\|\Phi_M\|_{L^2_G(M)}$: graph area
factors are bounded, and the stronger Gaussian absorbs the bounded
physical graph displacement, using the uniform bound for $Y^\perp$ on
$S_b$.  The annular height and slope are uniformly bounded.  From the
annular norm in Lemma~\ref{ref:real-comparison}, first choose
$\vartheta_{\rm cut}$ close enough to one and then enlarge $R_\eta$ to
obtain the bound $Ce^{-(1-\eta)R^2/8}$, assigning strictly smaller
losses to the cutoff fraction and the fixed additive radius change.
This is an $L^2$ norm estimate, which accounts for the denominator $8$.
The comparison lemma and \eqref{ref:taylor-multiplier} give
\begin{equation}\label{flow:comparison}
 \|w\|_{\mathcal H^2_{1+\delta}}+|\lambda_N|
      \le C(\mathcal D+|x|^{N+1}),\qquad
 |\nabla f_N(p)|\le C(\mathcal D+|x|^N).
\end{equation}
The finite-power inequality now yields
\[
 |x|^m\le C(\mathcal D+|x|^{N-1}).
\]
Shrinking the observation neighborhood absorbs the last term.
Consequently $|x|\le C\mathcal D^{1/m}$; if there are no opening
variables this step is omitted.  Our fixed choice of $N$ makes
the Taylor errors in \eqref{flow:comparison} and
\eqref{ref:taylor-action} at most $C\mathcal D^4$ when
$\mathcal D\le1$.  Hence \eqref{flow:comparison} is bounded by
$C\mathcal D$.  On $K_b$, where $\chi_R=1$ once $R\ge R_{\eta,K}$,
the identity $u=w+v_N$ gives the compact $H^2$ conclusion by adding
the finite Taylor height, whose compact norm is $O(|x|)$; take
$\gamma=\min\{1,1/m\}$, or $\gamma=1$ if $x$ is absent.
Finally smooth dependence of the finite multiplier jet gives
$|\lambda(b,0)-\lambda_N|\le C|x|$, proving
\eqref{flow:base-multiplier}.

Here is the action estimate, including the weight interpolation.
The finite-dimensional analytic gradient inequality (see
\cite[p.~1592]{Lojasiewicz1993} and
\cite[Section~0.2, equation~(0.4)]{CM2015}) for the single
analytic function $f_N$, at its critical point $(0,0)$, supplies
$\theta_0\in(0,1/2]$ with
\begin{equation}\label{flow:finite-dimensional-gradient}
 |f_N(b,x)-F(S)|^{1-\theta_0}
                     \le C|\nabla f_N(b,x)|.
\end{equation}
Choose $0<\theta<\min\{\theta_0,1/2\}$ and then choose
$\alpha<1$ close enough to one that
$2\alpha(1-\theta)>1$.  Take $\delta>0$ sufficiently small that
\[
 a_*:=\alpha(1+2\delta)+(1-\alpha)/2<1.
\]
For any of the finitely many polynomial factors in the second
variation, put $U=(1+|Y|)^k(|w|+|\nabla w|)$.  Polynomial
absorption and \eqref{flow:comparison} give
$\int U^2e^{-(1+2\delta)|Y|^2/4}\le C\mathcal D^2$.
The weak height and slope bounds, polynomial area growth, and a
\emph{weaker but still integrable} Gaussian give
$\int U^2e^{-|Y|^2/8}\le C$.  H\"older's inequality consequently
gives
\begin{equation}\label{flow:weighted-interpolation}
 \int U^2e^{-a_*|Y|^2/4}\,d\mu_{S_b}
 \le
 \left(\int U^2e^{-(1+2\delta)|Y|^2/4}\right)^\alpha
 \left(\int U^2e^{-|Y|^2/8}\right)^{1-\alpha}
 \le C\mathcal D^{2\alpha}.
\end{equation}
There is no use of finite unweighted area of $S_b$.

Integrate the first variation twice along the full graph segment
from $v_N$ to $v_N+w$.  Its first variation at $v_N$ consists of
the compact multiplier, of size $O(\mathcal D)$, paired with a
compact height of size $O(\mathcal D)$, and the high-order Taylor
residual.  The second variation of the area integrand involves
only $w,\nabla w$, with polynomial position factors; its Gaussian
on intermediate graphs is bounded by
$Ce^{-a_*|Y|^2/4}$ once tolerances are small.  Thus
\eqref{flow:weighted-interpolation} controls it.  The omitted
actual and graph tails are bounded, by \eqref{flow:gaussian-tail},
by $Ce^{-(1-\eta)R^2/4}\le C\mathcal D^2$; cutoff margins are
chosen with a strict exponent buffer here.  Together with
\eqref{ref:taylor-action}, this proves
\begin{equation}\label{flow:action-comparison}
 |F(M)-f_N(b,x)|\le C\mathcal D^{2\alpha}.
\end{equation}
Combine \eqref{flow:finite-dimensional-gradient},
\eqref{flow:comparison}, and \eqref{flow:action-comparison}.
Since $2\alpha(1-\theta)>1$ and $\theta<\theta_0$, the result is
\eqref{flow:localized-conclusions} for small $\mathcal D$.
For $\mathcal D$ bounded away from zero, it follows by increasing
the constant, using the area-ratio bound for $F(M)$ and the fixed
observation neighborhood for $x$.  All choices of $N$, the
analytic function, and $\theta,\gamma,\alpha,\delta$ preceded
the arbitrarily small localization loss $\eta$.  This proves
the claimed independence of the exponents.
\end{proof}

